\scanpage{53}
\mspage{52}

\uline{13 June} We have not yet expressed all the Axioms \uline{Cont}. 1, ...., there remain outstanding

\uline{Cont} 3 finite dimension

\uline{Cont} 4 ``admissibility'' -- but already follows I think from axioms C.

\uline{Cont} 6 \inserted{Criterion for} disjointness via the loci -- we shall return to this.

\uline{Cont} 7 Two loci with different dominant str. are disjoint.

But before continuing, I would like once again to change perspective !

Take as underlying set $\mathcal F$, the set of figures, with the two order relations \[ \leq\ \text{(or inclusion)} \qquad \text{and}\qquad \preccurlyeq\ \text{(subdivision)} \]

\uline{Country} (fourth version !)

\unsure{$U$ is|It is} a set $\mathcal F$ of objects, called \uline{figures}, and two order relations $\leq$ and $\preccurlyeq$ on $\mathcal F$ : \[ \mathcal C=(\underline{\mathcal F},\leq,\preccurlyeq) \] satisfying axioms\unsure{,|{}} which we are going to enumerate. If $F$ and $G$ are two figures and $F\leq G$, we say that $F$ is a \uuline{sub}-\uline{figure}, and we shall also often write $F\subset G$. If $F\preccurlyeq G$, we say that $F$ is a \uline{subdivision} of $G$, or that $F$ subdivides $G$.
